\documentclass[11pt]{article}
\usepackage[margin=1in]{geometry}
\usepackage[utf8]{inputenc}
\usepackage[T1]{fontenc}
\usepackage{amsmath, amsthm, amssymb, amsfonts}
\usepackage{hyperref}
\hypersetup{
  colorlinks=true,
  linkcolor=blue,
  urlcolor=cyan,
  pdftitle={The t-statistic is t-distributed},
}
\newtheorem{theorem}{Theorem}
\newtheorem{lemma}[theorem]{Lemma}
\newtheorem{proposition}[theorem]{Proposition}
\newtheorem{corollary}[theorem]{Corollary}
\theoremstyle{definition}
\newtheorem{definition}[theorem]{Definition}
\newtheorem{remark}[theorem]{Remark}
\title{The $t$-statistic is $t$-distributed}
\author{}
\date{\today}
\begin{document}
\maketitle
\begin{abstract}
We prove that the $t$-statistic of an OLS coefficient follows a Student
$t$-distribution with $n-p$ degrees of freedom, under the assumption of
Gaussian errors. The proof uses three tools: the sampling distribution
of $\hat\beta$, Cochran's theorem applied to the residual quadratic form,
and the independence of $\hat\beta$ and $\hat\sigma^2$.
\end{abstract}

\section{Introduction}
Richard showed up to Ethan's NBA Finals watch party with a notebook 
and pen in order to continue working on some proofs he wanted to write 
for work. I was inspired because I haven't written proofs since school, so
I thought I'd work on what he was working on Friday during slow work hours. 
It turns out, I misremembered, so this is a completely different proof to 
what he had.

When we fit a linear model and read off a coefficient $\hat\beta_j$, we want
to know whether the estimate is distinguishable from a hypothesized value
$\beta_j$ given the noise in the data. The $t$-statistic answers this by
standardizing the estimation error $\hat\beta_j - \beta_j$ by its estimated
standard deviation. The price of estimating $\sigma^2$ rather than knowing it
is that the resulting ratio is not normal but $t$-distributed, with heavier
tails.

The proof below generally follows dependency order. We first derive the sampling
distribution of $\hat\beta$ (Section~3.1), which gives a standard normal
numerator. We then take the residual-maker matrix $M$ (Section~3.2) and use
Cochran's theorem (Section~3.3) to show that the scaled residual sum of squares,
is $\chi^2_{n-p}$ (Section~3.4). After establishing that the
numerator and denominator are independent (Section~3.5), the main result
follows directly from the characterization of the $t$-distribution
(Section~4).

\section{Setup}

\subsection{The model and OLS estimator}
Given a regression problem $y = X\beta + \epsilon$ with unknown parameters $\beta$ and 
error $\epsilon \mid X \sim \mathcal{N}(0, \sigma^2 I)$, the 
ordinary least squares solution is $\hat\beta = (X^T X)^{-1} X^T y$.
Here $X$ is the $n \times p$ design matrix and $\beta$ is a $p \times 1$
vector. All distributional statements below are conditional on $X$.

\subsection{The $t$-statistic}
\begin{definition}
    The $t$-statistic for coefficient $\beta_j$ is defined as
    $$t_j = \frac{\hat\beta_j - \beta_j}{\operatorname{S.E.}(\hat\beta_j)}
          = \frac{\hat\beta_j - \beta_j}{\sqrt{\hat\sigma^2 \left[(X^T X)^{-1}\right]_{jj}}}$$
\end{definition}
\subsection{The error assumptions}
For $n$ observations, the error terms $\epsilon$ are an $n \times 1$ vector satisfying the conditions 
$\mathbb{E}(\epsilon) = 0$ and $\operatorname{Var}(\epsilon) = \sigma^2 I_n$, where $\sigma^2$ is the error variance.

Our assumption is that a fitted model has uncorrelated and constant variance across all observations.

Standard error is the estimated standard deviation of a statistic. Since
$\sigma^2$ is unknown, we must estimate it from the residuals; the estimator
$\hat\sigma^2$ is constructed in Section~3.4.

\section{Building blocks}

\subsection{Sampling distribution of $\hat\beta$}
To find the standard error, we first compute $\operatorname{Var}(\hat\beta \mid X)$ and substitute the population variance $\sigma^2$
for the sample variance $\hat\sigma^2$.
\begin{lemma}\label{lem:beta-dist}
    $\hat\beta \,|\, X \sim \mathcal{N}(\beta, \sigma^2 (X^T X)^{-1})$
\end{lemma}

\begin{proof}
    Substituting $y = X\beta + \epsilon$,
    $$\hat\beta = (X^T X)^{-1} X^T (X\beta + \epsilon) = \beta + (X^T X)^{-1} X^T \epsilon.$$
    So $\hat\beta$ is an affine transformation of the Gaussian vector
    $\epsilon$ and is therefore Gaussian. Its mean is
    $$\mathbb{E}(\hat\beta \mid X) = \beta + (X^T X)^{-1} X^T \, \mathbb{E}(\epsilon \mid X) = \beta.$$
    For the variance, let $A = (X^T X)^{-1} X^T$ and recall $\operatorname{Var}(\epsilon \mid X) = \sigma^2 I$. Then,
\begin{align*}
    \operatorname{Var}(\hat\beta \mid X) &= \operatorname{Var}(A y \mid X) \\
                                &= A \operatorname{Var}(y \mid X) A^T \\
                                &= (X^T X)^{-1} X^T (\sigma^2 I) \big[ (X^T X)^{-1} X^T\big]^T \\
                                &= (X^T X)^{-1} X^T (\sigma^2 I) X (X^T X)^{-1} \\
                                &= \sigma^2 (X^T X)^{-1} \qedhere
\end{align*}
\end{proof}

Thus, we have one building block for the $t$-distribution: for each
coordinate $j$,

$$ \hat\beta_j - \beta_j \sim \mathcal{N}\big(0, \, \sigma^2 \left[(X^T X)^{-1}\right]_{jj}\big) \implies \frac{\hat\beta_j - \beta_j}{\sqrt{\sigma^2 \left[(X^T X)^{-1}\right]_{jj}}} \sim \mathcal{N}(0, 1) $$

And, the standard error

$$\operatorname{S.E.}(\hat\beta_j) = \sqrt{\hat\sigma^2 \left[(X^T X)^{-1}\right]_{jj}}.$$

In order to develop a definition for $\hat\sigma^2$, we need to first define the residuals $\hat\epsilon$ or $e$.

\begin{align*}
    e &= y - X\hat\beta \\
      &= y - X (X^T X)^{-1} X^T y \\
      &= (I - X (X^T X)^{-1} X^T) y
\end{align*}

\subsection{The residual-maker matrix or annihilator matrix}
Since we are estimating the linear model with OLS, our fitted values are 
$$\hat y = X \hat\beta = X (X^T X)^{-1} X^T y = \mathbf{P} y$$
Thus, our projection matrix is 
$$\mathbf{P} = X(X^T X)^{-1} X^T.$$
\begin{lemma}\label{lem:residual-maker}
    $M = I - \mathbf{P} = I - X(X^T X)^{-1} X^T$ is symmetric, idempotent, has rank $n-p$, satisfies $MX= 0$, and $e = M\epsilon$.
    $M$ is known as the residual-maker, or the annihilator matrix. 
\end{lemma}

\begin{proof}
    \emph{Symmetry.} Since $(X^T X)^{-1}$ is symmetric,
    $$M^T = I - \big[X (X^T X)^{-1} X^T\big]^T = I - X (X^T X)^{-1} X^T = M.$$

    \emph{Idempotency.} First, $\mathbf{P}$ is idempotent:
    $$\mathbf{P}^2 = X (X^T X)^{-1} \underbrace{X^T X (X^T X)^{-1}}_{= \, I} X^T = X (X^T X)^{-1} X^T = \mathbf{P}.$$
    Therefore
    $$M^2 = (I - \mathbf{P})^2 = I - 2\mathbf{P} + \mathbf{P}^2 = I - \mathbf{P} = M.$$

    \emph{Annihilation.} 
    $$MX = X - X (X^T X)^{-1} X^T X = X - X = 0.$$

    \emph{Rank.} For an idempotent matrix, rank equals trace. Using the
    cyclic property of the trace,
    $$\operatorname{rank}(M) = \operatorname{tr}(M) = \operatorname{tr}(I_n) - \operatorname{tr}\big(X (X^T X)^{-1} X^T\big) = n - \operatorname{tr}\big((X^T X)^{-1} X^T X\big) = n - \operatorname{tr}(I_p) = n - p.$$

    \emph{Residuals.} Substituting $y = X\beta + \epsilon$ and using $MX = 0$,
    $$e = My = M(X\beta + \epsilon) = MX\beta + M\epsilon = M\epsilon. \qedhere$$
\end{proof}


\subsection{Cochran's theorem}
\begin{theorem}[Cochran's Theorem]\label{thm:cochran}
    Let $z \sim \mathcal{N}(0, I_n)$ and suppose
    $$z^T z = \sum^m_{j=1} Q_j, \qquad Q_j = z^T A_j z,$$
    where each $A_j$ is symmetric and $\sum^m_{j=1} \operatorname{rank}(A_j) = n$.
    Then the $Q_j$ are mutually independent and
    $Q_j \sim \chi^2_{\operatorname{rank}(A_j)}$ for each $j$.
\end{theorem}

\begin{proof}
    Sorry, we omit the proof, so see Keener, \emph{Theoretical Statistics} (2010).
\end{proof}

In our application we take $z^T z = z^T \mathbf{P} z + z^T (I - \mathbf{P}) z$,
with $\operatorname{rank}(\mathbf{P}) + \operatorname{rank}(I - \mathbf{P}) = p + (n - p) = n$.

\subsection{Distribution of $\hat\sigma^2$}

\begin{definition}[Residual sum of squares]
    $\mathrm{RSS} = e^T e = \left\| (I - \mathbf{P}) y \right\|^2$.
\end{definition}

\begin{definition}[Variance estimator]
    $\displaystyle \hat\sigma^2 = \frac{\mathrm{RSS}}{n - p}$.
\end{definition}

We express $\mathrm{RSS}$ as a quadratic form in the errors $\epsilon$.

\begin{lemma}\label{lem:rss-quadform}
    $\mathrm{RSS} = \epsilon^T (I - \mathbf{P})\, \epsilon$.
\end{lemma}

\begin{proof}
    By Lemma~\ref{lem:residual-maker}, $e = M\epsilon = (I - \mathbf{P})\,\epsilon$. Since $M$ is symmetric ($M^T = M$) and idempotent
    ($M^2 = M$), 
    \begin{align*}
        \mathrm{RSS} &= e^T e \\
                     &= (M\epsilon)^T (M\epsilon) \\
                     &= \epsilon^T M^T M \epsilon \\
                     &= \epsilon^T M^2 \epsilon \\
                     &= \epsilon^T M \epsilon \qedhere
    \end{align*}
\end{proof}

\begin{proposition}\label{prop:rss-chisq}
    $\displaystyle \frac{\mathrm{RSS}}{\sigma^2} \sim \chi^2_{n-p}$.
\end{proposition}

\begin{proof}
    Because Cochran's theorem applies to standard normal vectors, we scale 
    $\epsilon \sim \mathcal{N}(0, \sigma^2 I)$ by $\epsilon / \sigma = z \sim \mathcal{N}(0, I_n)$.
    With Lemma~\ref{lem:rss-quadform},
    $$
        z^T (I - \mathbf{P})\, z
        = \frac{\epsilon^T (I - \mathbf{P})\, \epsilon}{\sigma^2}
        = \frac{\mathrm{RSS}}{\sigma^2}.
    $$
    By Lemma~\ref{lem:residual-maker}, $\operatorname{rank}(I - \mathbf{P}) = n - p$, so 
    Cochran's theorem gives us that the quadratic form $z^T (I - \mathbf{P})\, z$
    follows a $\chi^2$ distribution with $n - p$ degrees of freedom.
\end{proof}

\begin{corollary}\label{cor:s2-dist}
    $\displaystyle \frac{(n-p)\,\hat\sigma^2}{\sigma^2} \sim \chi^2_{n-p}$.
\end{corollary}

\begin{proof}
    Follows from Proposition~\ref{prop:rss-chisq}, since
    $(n-p)\,\hat\sigma^2 = \mathrm{RSS}$.
\end{proof}

\begin{corollary}[Unbiasedness]
    $\mathbb{E}\big(\hat\sigma^2\big) = \sigma^2$.
\end{corollary}

\begin{proof}
    A $\chi^2_{n-p}$ random variable has mean $n-p$. Taking expectations in
    Corollary~\ref{cor:s2-dist},
    $\mathbb{E}\big((n-p)\,\hat\sigma^2/\sigma^2\big) = n-p$, and rearranging
    gives $\mathbb{E}\big(\hat\sigma^2\big) = \sigma^2$.
\end{proof}

\subsection{Independence of $\hat\beta$ and $\hat\sigma^2$}
The last building block is that the numerator and denominator of the
$t$-statistic do not share information.

\begin{lemma}\label{lem:independence}
    Conditional on $X$, $\hat\beta$ and $\hat\sigma^2$ are independent.
\end{lemma}

\begin{proof}
    From the proof of Lemma~\ref{lem:beta-dist}, $\hat\beta - \beta = A\epsilon$
    with $A = (X^T X)^{-1} X^T$, and from Lemma~\ref{lem:residual-maker},
    $e = M\epsilon$. Both are linear in the same Gaussian vector $\epsilon$,
    so $(\hat\beta - \beta, \, e)$ is jointly Gaussian. Their cross-covariance is
    \begin{align*}
        \operatorname{Cov}(A\epsilon, M\epsilon) &= A \operatorname{Var}(\epsilon) M^T \\
                                                 &= \sigma^2 A M \\
                                                 &= \sigma^2 (X^T X)^{-1} X^T M \\
                                                 &= \sigma^2 (X^T X)^{-1} (M X)^T \\
                                                 &= 0
    \end{align*}
    using symmetry of $M$ and $MX = 0$. Uncorrelated jointly Gaussian vectors
    are independent, so $\hat\beta$ is independent of $e$. Since
    $\hat\sigma^2 = e^T e / (n-p)$ is a function of $e$ alone, $\hat\beta$ is
    independent of $\hat\sigma^2$.
\end{proof}

\section{The Student $t$-distribution}
\begin{definition}\label{def:t-dist}
    Let $Z \sim \mathcal{N}(0,1)$ and $W \sim \chi^2_m$ be independent. Then
    $$t = \frac{Z}{\sqrt{W / m}}$$
    follows the Student $t$-distribution with $m$ degrees of freedom.
\end{definition}

\section{Main result}
\begin{theorem}
    Under the model of Section~2, $t_j \sim t_{n-p}$.
\end{theorem}

\begin{proof}
    Write $c_{jj} = \left[(X^T X)^{-1}\right]_{jj}$. Dividing the numerator and
    denominator of $t_j$ by $\sqrt{\sigma^2 c_{jj}}$,
    \begin{align*}
        t_j &= \frac{\hat\beta_j - \beta_j}{\sqrt{\hat\sigma^2 \, c_{jj}}} \\
            &= \frac{\big(\hat\beta_j - \beta_j\big) \big/ \sqrt{\sigma^2 c_{jj}}}{\sqrt{\hat\sigma^2 / \sigma^2}} \\
            &= \frac{Z}{\sqrt{W / (n-p)}},
    \end{align*}
    where
    $$Z = \frac{\hat\beta_j - \beta_j}{\sqrt{\sigma^2 c_{jj}}} \sim \mathcal{N}(0,1)
    \qquad \text{and} \qquad
    W = \frac{(n-p)\,\hat\sigma^2}{\sigma^2} \sim \chi^2_{n-p}$$
    by Lemma~\ref{lem:beta-dist} and Corollary~\ref{cor:s2-dist} respectively.
    See that the unknown $\sigma^2$ cancels in the ratio. By
    Lemma~\ref{lem:independence}, $Z$ (a function of $\hat\beta$) and $W$ (a
    function of $\hat\sigma^2$) are independent. By
    Definition~\ref{def:t-dist}, $t_j \sim t_{n-p}$.
\end{proof}

\section{Remarks}
\begin{remark}
    Statistical libraries test the null hypothesis $H_0 \colon \beta_j = 0$,
    where the numerator then reduces to $\hat\beta_j$ and the reported
    statistic is $t_j = \hat\beta_j / \operatorname{S.E.}(\hat\beta_j)$. Thus, a
    large $|t_j|$ relative to the $t_{n-p}$ distribution indicates that the
    coefficient is statistically significant.
\end{remark}

\end{document}

%\begin{remark}
%    The normality assumption $\epsilon \mid X \sim \mathcal{N}(0, \sigma^2 I)$
%    is what makes the result exact in finite samples. Without it, $t_j$ is
%    only asymptotically normal, by the central limit theorem.
%\end{remark}


%\begin{corollary}[Unbiasedness]
%    $\mathbb{E}\big[\hat\sigma^2\big] = \sigma^2$.
%\end{corollary}
%
%\begin{proof}
%    A $\chi^2_{n-p}$ random variable has mean $n - p$. Taking expectations
%    in Corollary~\ref{cor:s2-dist},
%    $\mathbb{E}\big[(n-p)\,\hat\sigma^2 / \sigma^2\big] = n - p$,
%    and rearranging gives $\mathbb{E}\big[\hat\sigma^2\big] = \sigma^2$.
%\end{proof}

